%==============================================================================
% MATH 5080 — Introduction to Analysis I (Fall 2026)
% Class:        7
% Date:         Sept 15th
%==============================================================================

\begin{theorem}[Operation on limits]
    Suppose 
    $\lim_{n \to \infty} a_n = A$,
    $\lim_{n \to \infty} b_n = B$, then 
    \begin{enumerate}
        \item $\lim_{n \to \infty}(a_n + b_n) = A + B$
        \item $\lim_{n \to \infty}ka_n = kA$
        \item $\lim_{n \to \infty}a_n b_n = AB$
        \item If $b_n \neq 0, \forall n$ and $B \neq 0$, then 
        \[
            \lim_{n \to \infty} \frac{a_n}{b_n} = \frac{A}{B}
        \]
    \end{enumerate}
\end{theorem}

\begin{proof}[Proof of 3]
    \textit{
        The idea is to use boundedness of one of the sequences to 
        find a bound on the deviation for each of the convergent sequences
        from their limits, and then take the bigger $N$.
    }
    \medskip

    Fix $\ep > 0$, seek $N > 0$ s.t. $\forall n > N, \left| a_n b_n - AB \right| < \ep $.
    \begin{align*}
        \left| a_n b_{n} - AB \right|  
        &= \left| a_{b}b_{n} - a_{b}B + a_{b}B - AB \right| \\
        &\leq \left| a_{n}b_{n} - a_{n}B \right| + \left| a_{n}B - AB \right| \\
        &= \left| a_{n} \right|  |b_{b} - B| + |a_{n} - A | |B|
    \end{align*}
    Since convergent implies boundedness, $\exists M > 0$ s.t. $| a_n | \leq M, \forall n$.
    Choose $M$ such that $|B| \leq M$. 
    \medskip

    By convergence of $a_n, b_n$, we can choose
    \begin{itemize}
        \item $N_1 > 0$ s.t. $|a_n - A | < \frac{\ep}{2M}, \forall n > N_1$
        \item $N_2 > 0$ s.t. $|a_n - A | < \frac{\ep}{2M}, \forall n > N_2$
        \item $N = \max \{ N_1, N_2 \} $
    \end{itemize} 
    Then $\forall n > N$
    \begin{align*}
        \left| a_n b_{n} - AB \right|  
        &\leq \left| a_{n} \right|  |b_{b} - B| + |a_{n} - A | |B| \\
        &< |a_n | \cdot \frac{\ep}{2M} + \frac{\ep}{2M} \cdot |B| \\
        & \leq \frac{\ep}{2} + \frac{\ep}{2} \\
        &= \ep
    \end{align*}
\end{proof}

\begin{example}
    $\lim_{n \to \infty} \frac{n + 1}{2n - 1} = \frac{1}{2}$ 
    \medskip

    By definition, fix $\ep > 0$, seek $N > 0$ s.t. $\forall n > N$, 
    \[
        \left| 
            \frac{n+1}{2n-1} - \frac{1}{2}
        \right|  < \ep
    \]
    Simplifying the expression, we get
    \begin{align*}
        & \left| 
            \frac{n+1}{2n-1} - \frac{1}{2}
        \right|  
        = \left| \frac{2n + 2 - 2n +1}{2(2-n-1)} \right| 
        = \frac{3}{2} \frac{1}{2n - 1} < \ep \\
        \iff &  \frac{1}{2n - 1} < \frac{2}{3}\ep \\
        \iff & 2n - 1 > \frac{3}{2} \ep \\
        \iff & n > \frac{1}{2} \left( \frac{3}{2} \ep  + 1 \right) 
    \end{align*}
    Take 
    \[
        N= 
        \lceil
        \frac{1}{2} \left( \frac{3}{2} \ep + 1 \right)
        \rceil
    \]
    and we are done. 

    \medskip

    Alternatively, \textit{we can prove using algebraic properties}.
    \begin{align*}
        \frac{n + 1}{2n - 1} 
        &= \frac{ \frac{1}{2} (2n-1) + \frac{3}{2}}{2n - 1} \\
        &= \frac{1}{2} + \frac{3}{2} \frac{1}{2n -1}
    \end{align*}
    It suffices to prove $\lim_{n  \to \infty} \frac{1}{2n-1} = 0$.
    Then by theorem, $\lim_{n \to \infty} \frac{n + 1}{2n-1} = \frac{1}{2} + \frac{3}{2} \cdot 0 = \frac{1}{2}$.
\end{example}


\begin{corollary}
    If $\lim_{n \to \infty}a_n$ converges, and $\lim_{n \to \infty} b_n$ diverges,
    then $\lim_{n \to \infty}a_n + b_n$ diverges.
\end{corollary}

\begin{example}
    Show $\lim_{n \to \infty} a^{ \frac{1}{n}} = 1$ 
    \medskip

    \textbf{Proof}
    We consider 3 cases
    \begin{itemize}
        \item $a < 1$: then by theorem,
        \begin{align*}
            &\lim_{n \to \infty} \frac{1}{a^{ \frac{1}{n}}}
            = \frac{1}{ \lim_{n \to \infty} a^{ \frac{1}{n}}} \\
            \implies & \lim_{n \to \infty} a^{ \frac{1}{n}} = 
            \frac{1}{\lim_{n \to \infty} \frac{1}{a^{ \frac{1}{n}}}}
            = \frac{1}{\lim_{n \to \infty} \left( \frac{1}{a} \right)^{ \frac{1}{n}} }
        \end{align*}
        and then we are in the case of $a > 1$
        \item $a = 1$: trivial
        \item $a > 1$
    \end{itemize} 

    Hence it suffices to show for $a > 1$. 
    Fix $\ep > 0$, we seek $N$ s.t. $\forall n > N$,
    \begin{align*}
        & \left| a^{ \frac{1}{n}} - 1 \right| < \ep \\
        \iff & a^{ \frac{1}{n}} - 1 < \ep \\
        \iff & a^{ \frac{1}{n}}  < 1 + \ep \\
        \iff & \frac{1}{n} \ln a < \ln( 1 + \ep) \\
        \iff & n > \frac{ \ln a}{\ln ( 1 + \ep)}
    \end{align*}
    Hence we can take 
    \[
        N = 
        \lceil
        \frac{\ln a}{\ln ( 1 + \ep)}
        \rceil
    \]
    and we are done.
\end{example}


\begin{theorem}[Sandwich Theorem]
    Let $N_0$ such that $a_n \leq c_n \leq b_n$,
    $\forall n > N_0$,
    suppose
    \[
        \lim_{n \to \infty} a_n =
        \lim_{n \to \infty} b_n = L
    \]
    then 
    \[
        \lim_{n \to \infty} c_n = L
    \]
\end{theorem}

\begin{proof}
    Fix $\varepsilon > 0$, we seek $N > 0$ s.t. $\forall n > N$,
    \[
        |c_n - L| < \varepsilon
    \]
    Since $\lim a_n = L, \exists N_1$ s.t. 
    $\forall n > N_1$,
    $|a_n - L | < \ep$
    Since $\lim b_n = L, \exists N_2$ s.t. 
    $\forall n > N_2$,
    $|b_n - L | < \ep$
    Pick $N = \max \{ N_1, N_2, N_0 \} $, 
    then $\forall n > N$,
    \begin{align*}
        c_n - L \leq b_n - L < \ep \\
        c_n - L \geq a_n - L > -\ep \\
    \end{align*}
    Hence 
    \[
        \left| c_n - L  \right|  < \ep
    \]
    
\end{proof}

\begin{example}
    Show $\lim_{n \to \infty} n^{ \frac{1}{n}} = 1$ 
    \medskip

    By sandwich theorem, we first have $n^{ \frac{1}{n}} \geq 1$
    for all $n \geq 1$. Hence we have found a lower bound 
    for all terms with $n > 1$.
    \medskip

    We then seek an upper bound that goes to 1. 
    Define a helper $h_n = n^{ \frac{1}{n}} - 1$,
    then $n^{ \frac{1}{n}} = 1 + h_n$. 
    We want to find an upper bound of $h_n$ that 
    goes to 0. 
    Write 
    \begin{align*}
        n^{ \frac{1}{n}} & = 1 + h_n \\
        \left( n^{\frac{1}{n}} \right)^{n} & = (1 + h_n)^n \\
        &= \sum\limits_{k = 0}^{n} 
        \begin{pmatrix} n \\ k \end{pmatrix} 
        h_n^k \\
        & \geq \begin{pmatrix} n \\2 \end{pmatrix} h_n^2 \\
        &= \frac{n(n-1)}{2} h_n^2 \\
        \frac{2n}{n(n-1)} &\geq h_n^2 \\
        \sqrt{ \frac{2}{n-1}} & \geq h_n
    \end{align*}
    Hence we can apply the sandwich theorem and say that 
    $n^{ \frac{1}{n}} \to 1$.
\end{example}

\begin{definition}[Divergence]
    We say $\{ a_n \} $ diverges to $\infty$ if $\forall M > 0$,
    $\exists N > 0$ s.t. $a_n > M, \forall n > N$, denoted 
    \[
        \lim_{n \to \infty} a_n = \infty
    \]
\end{definition}

\begin{example}
    $\{ 0, 2, 0, 4, \hdots \} $ does not diverge to $\infty$.
\end{example}

\subsection{Cauchy sequences and monotone sequences}

\begin{definition}[Cauchy sequence]
    $\{ a_n \} $ is a \textbf{Cauchy sequence} if $\forall \varepsilon > 0$, 
    $\exists N > 0$ s.t. $|a_n - a_m| < \varepsilon, \forall n, m > N$.
\end{definition}

\begin{theorem}[Cauchy implies boundedness]
    Every Cauchy sequence is bounded.
\end{theorem}
\begin{proof}
    Fix $\{ a_n \} $ a Cauchy sequence. 
    Fix $\varepsilon = 1$, by definition, 
    $\exists N_0 > 1$ s.t. $\forall n, m > N_0$,
    $|a_n - a_m| < 1$. 
    Hence, $a_n < |a_{N_0 + 1}| + 1, \forall n > N_0$. 
    \medskip 

    Also, $\{ a_n \}_{1 \leq n \leq N_0} $ is finite 
    and a maximum exists. 
    \medskip

    Hence $M \coloneq \max \{ 
        a_1, \hdots a_{N_0},
        |a_{N_0 + 1}|
    \} 
    $ bounds the whole sequence.
\end{proof}

\begin{theorem}[Convergence imply Cauchyness]
    Every convergent sequence is Cauchy.
\end{theorem}
\begin{proof}
    Fix sequence $ \{ a_n \} $, fix $ \varepsilon > 0$. Let limit be $L$. 
    Since $a_n$ converges, $\exists N > 0$ s.t. 
    \begin{align*}
        & |a_n - L | < \frac{\varepsilon}{2}, \forall n > N\\
        \implies & \forall n, m > N, |a_n - L| < \frac{\varepsilon}{2}, 
        |a_m - L | < \frac{\varepsilon}{2} \\
        \implies & |a_n - a_m| = 
        |(a_n - L) - (a_m - L)| \leq |a_n - L | + |a_m - L | < 
        \frac{\varepsilon}{2} + 
        \frac{\varepsilon}{2} = \varepsilon
    \end{align*}
\end{proof}


\begin{theorem}[Cauchy Convergence Criterion]
    If $\{ a_n \} \in \mathbb{R}$ is Cauchy then it 
    converges, and the limit is in $\mathbb{R}$. 
\end{theorem}

\begin{proof}
    \textit{We use Bolzano-Weierstrass to show existence of limit.}
    \medskip 

    Consider two cases. $\{ a_n \} $ is either finite or infinite. 
    \begin{itemize}
        \item $\{ a_n \} $ finite. Let $S$ dnote the set of values $ \{ s_1, \hdots s_k \} $. 
        Let $\varepsilon = \frac{1}{2} \min \{ |s_i - s_j| \}, i \neq j, 1 \leq i, j \leq k $. 
        \medskip 

        Since $\{ a_n \} $ is Cauchy, $\exists N > 0$ s.t. 
        \[
            |a_n - a_m| < \varepsilon, \forall n, m > N \implies 
            \exists i \text{ s.t. } a_n = s_i, \forall n > N
        \]
        \item $ \{ a_n \} $ infinite. 
        We have shown above that Cauchyness implies boundedness.
        By BW, $\exists L \in \mathbb{R}$
        that is an accumulation point. 
        \medskip 

        WTS $\lim_{n \to \infty}a_n = L$. Fix $\varepsilon > 0$, 
        we seek $N$ s.t. $|a_n - L| < \varepsilon, \forall n > N$. 
        \medskip 

        Since $\{ a_n \} $ Cauchy, $\exists N_0$ s.t. 
        $|a_n - a_m | < \frac{ \varepsilon}{2}, \forall n, m > N$.
        If suffices to show that at least one
        $a_{n_0} \in N\left(L, \frac{\varepsilon}{2} \right), n_0 > N_0$.

        Since $L$ is an accumulation point, $N^*\left(L, \frac{\varepsilon}{2}\right)$
        contains infinitely many points, so in particular we may pick
        $n_0 > N_0$ with $a_{n_0} \in N\left(L, \frac{\varepsilon}{2}\right)$.
        Then for $n > n_0$,
        \[
            |a_n - L| \leq |a_n - a_{n_0}| + |a_{n_0} - L|
            < \frac{\varepsilon}{2} + \frac{\varepsilon}{2} = \varepsilon,
        \]
        i.e. $N\left(a_{n_0}, \frac{\varepsilon}{2}\right) \subseteq N(L,\varepsilon)$
        and every later term lies in it.
        \medskip

        \begin{center}
        \begin{tikzpicture}[scale=.80]
            % Outer ball N(L, epsilon)
            \draw[thick, blue] (0,0) circle (3);
            \node[blue] at (-2.1,2.1) {$N(L,\varepsilon)$};

            % Inner ball N(L, epsilon/2)
            \draw[thick, dashed, gray] (0,0) circle (1.5);
            \node[gray] at (0.6,-1.75) {$N(L,\varepsilon/2)$};

            \fill (0,0) circle (1.5pt) node[below left] {$L$};

            % Point a_{n_0} inside the inner ball
            \coordinate (an0) at (0.6,0.6);

            % Ball N(a_{n_0}, epsilon/2)
            \draw[thick, red] (an0) circle (1.5);
            \node[red] at ($(an0)+(0,1.75)$) {$N(a_{n_0},\varepsilon/2)$};
            \fill (an0) circle (1.5pt) node[above right] {$a_{n_0}$};

            % Later terms a_n, n > n_0, inside N(a_{n_0}, epsilon/2)
            \coordinate (ann) at ($(an0)+(0.9,-0.4)$);
            % \fill ($(an0)+(-0.5,0.8)$) circle (1pt);
            % \fill ($(an0)+(0.3,1.0)$) circle (1pt);
            % \fill ($(an0)+(-0.9,-0.5)$) circle (1pt);

            % Triangle inequality: L -- a_{n_0} -- a_n -- L
            \draw[thick, orange!80!black] (0,0) -- (an0) -- (ann) -- cycle;
            \node[orange!80!black] at ($(0,0)!0.5!(an0) + (0.3,0.05)$) {\scriptsize $\varepsilon/2$};
            \node[orange!80!black] at ($(an0)!0.5!(ann) + (0.35,0.05)$) {\scriptsize $\varepsilon/2$};
            \node[orange!80!black] at ($(0,0)!0.5!(ann) + (-0.3,-0.25)$) {\scriptsize $<\varepsilon$};
            \fill (ann) circle (1.5pt) node[below right] {\scriptsize $a_n,\, n>n_0$};
        \end{tikzpicture}
        \end{center}
    \end{itemize}
\end{proof}

\begin{example}
    We can show divergence of Harmonic series by showing that 
    it is not Cauchy. 
    \[
        a_n = 1 + \frac{1}{2} + \hdots + \frac{1}{n}
    \]
    i.e. WTS $\exists \ep >0$,
    $\forall N > 0$,
    $\exists m > n > N$ s.t. 
    $|a_m - a_n |> \ep$.
    Take $\ep  = \frac{1}{4}$, 
    $m = 2n$,
    then 
    \begin{align*}
        a_m- a_n &= 
        \frac{1}{ n+1} + 
        \frac{1}{ n+2} + 
        \hdots
        \frac{1}{m}
        \\
        & \geq \frac{m - n}{m} \\
        & = 1 - \frac{n}{m} \\
        & = \frac{1}{2}
    \end{align*}
    Hence the harmonic series is not cauchy and therefore does not converge.
\end{example}


\begin{definition}[Monotone sequence]
    \hfill
    \begin{enumerate}
        \item $\{ a_n \} $ is \textbf{increasing} if $a_n \leq a_{n + 1}, \forall n$
        \item $\{ a_n \} $ is \textbf{decreasing} if $a_n \geq a_{n + 1}, \forall n$
    \end{enumerate}
    If $ \{ a_n \} $ is either increasing or 
    decreasing, we say it is \textbf{monotone}.
    \medskip

    Equivalently, we can define strictly increasing, strictly decreasing and
    strictly monotone.
\end{definition}

\begin{theorem}[Monotone Convergence Theorem]
    If a sequence $\{ a_n \} $ is monotone and bounded,
    then it converges. 
\end{theorem}

\begin{corollary}[Alternative statement of MCT]
    Given a monotone sequence, it is bounded if and only
    if it converges.
\end{corollary}

\begin{proof}[Proof of MCT]
    WTS that given a monotone sequence, 
    it converges if and only it is bounded.
    \medskip

    $(\implies)$: convergence implies boundedness always, 
    without requiring monotonicity.
    \medskip

    $(\impliedby)$:
    Let $\{ a_n \} $ be a monotone and bounded sequence, 
    WTS it converges. 
    \medskip

    WLOG, assume $\{ a_n \} $ increasing. 
    Since $ \{ a_n \} $ bounded, by CA, 
    $\sup \{ a_n \} $ exists, say $L$. 
    WTS $\lim_{n \to \infty}a_n = L$.
    \medskip 

    Fix $\varepsilon > 0$< by definition of supremum, 
    $\exists n_0$ s.t. $a_{n_0} > L-\varepsilon$.
    Hence 
    \begin{align*}
        & L - \varepsilon < a_n \leq L, \forall n > n_0 \\
        \implies & |a_n - L | < \varepsilon, \forall n > n_0
    \end{align*}
\end{proof}

\begin{remark}
    \hfill
    \begin{enumerate}
        \item An increasing sequence converges $\iff$ bounded from above.
        \item A decreasing sequence converges $\iff$ bounded from below.
    \end{enumerate}
\end{remark}

\begin{example}
    \textit{We can use MCT to determine the convergence
    of a sequence defined recursively,
    i.e. we prove convergence by showing
    boundedness and monotonicitiy}
    \medskip

    For example, $a_1 = 5$, $a_{n + 1} = \frac{1}{2} \left( a_n + \frac{5}{a_n} \right) $.
    We can show this is bounded by noting $a_n > 0$ and 
    therefore 
    \begin{align*}
        a_{n + 1} &= \frac{1}{2} \left( 
            \sqrt{a_n}^2 + \sqrt{ \frac{5}{a_n}}^2
        \right)  \\
        & \geq \sqrt{a_n} \sqrt{ \frac{5}{a_n}}
        = \sqrt{5}
    \end{align*}
\end{example}



\begin{proof}[Proof that interval are connected 
    \label{proof:connected-sets-in-R}]
    \medskip

    Recall theorem \ref{theorem:connected-sets-in-R}. 
    WTS that if $S \subset \mathbb{R}$ is an interval
    then $S$ is connected.
    \medskip

    Suppose $S$ disconnected. By definition, $\exists A, B$
    open and nonempty such that $S \subset A \cup B$,
    $A \cap B = \emptyset$,
    $S \cap A \neq \emptyset$,
    $S \cap B \neq \emptyset$.
    \medskip

    Let $x \in S \cap A, y \in S \cap B$. WLOG, say $x < y$.
    Let $x_1 = x, y_1 = y$. We seek sequences 
    $\{ x_n \} \subset A \cap S $,
    $\{ y_n \} \subset A \cap S $,
    such that 
    \begin{itemize}
        \item $x_n$ is increasing 
        \item $y_n$ is decreasing 
        \item $|x_n - y_n| \to 0$ as $n \to \infty$
    \end{itemize} 

    To construct such a sequence, take $z_{k + 1} = \frac{x_k + y_k}{2}$.
    \begin{itemize}
        \item if $z_{k + 1} \in A, x_{k + 1} = z_{k + 1}, y_{k + 1} = y_k$
        \item if $z_{k + 1} \in B, x_{k + 1} = x_k, y_{k + 1} = z_{k + 1}$
    \end{itemize} 
    Then 
    \[
        y_i - x_i = \frac{y_1 - x_1}{2^{i - 1}}
    \]
    Since $\{ x_n \} $ bounded and increasing, 
    $\exists a \in \mathbb{R}$ s.t. $x_n \to a$.
    Likewise
    $\exists b \in \mathbb{R}$ s.t. $y_n \to b$.
    Then
    \[
        b - a = 
        \lim_{n \to \infty} y_n - 
        \lim_{n \to \infty} x_n 
        = \lim_{n  \to \infty} (y_{n} - x_{n})
        = \lim_{n  \to \infty} 2^{-(n-1)}(y_1 = x_1)
        = 0
    \]
    Hence $b=a$.
    \medskip

    We know that $a \in (S \cap A) \setminus ( S \cap B)$ or 
    $a \in (S \cap B) \setminus (S \cap A)$. 
    Suppose $a \in S \cap A$ but $a \notin B$, by openess of $A$, 
    there is some $N(a, \varepsilon) \subset A$.
    Since $A, B$ disjoint, $N(a, \varepsilon) \subset B^c$. 
    However, $y_n \to a$, hence $\exists N>0$ s.t. $\forall n > N$,
    $y_n \in N(a, \varepsilon)$, which is a contradiction.
\end{proof}

\subsection{Subsequences}

\begin{definition}[Subsequence]
    Let $\{ a_n \}  \in \mathbb{R}$ and $\{ n_k \}_k$ be a 
    striclty increasing sequence 
    in $\mathbb{N}$, i.e. $n_k < n_{k + 1}, \forall k$, then 
    $\{ a_{n_k} \}_k$ is a subsequence of $\{ a_n \} $
\end{definition}

\begin{example}
    Take $\{ 1, -1, 1, -1, \hdots \} $, 
    \begin{itemize}
        \item $n_k = 2k - 1, \{ n_k \}_k = \{ 1, 3, 5, \hdots \}$, 
        $\{ a_{n_k} \}_k = \{ 1, 1, 1, \hdots \}  $
        \item $n_k = 2k , \{ n_k \}_k = \{ 2, 4, 6, \hdots \}$, 
        $\{ a_{n_k} \}_k = \{ -1, -1, -1, \hdots \}  $
    \end{itemize} 
\end{example}

\begin{theorem}[Subsequence of convergent sequence converges]
    If $\lim_{n \to \infty}a_n = L \in \mathbb{R}$, or $L = \pm \infty$
    then for any subsequence $\{ a_{n_k} \}_k $ of $\{ a_n \} $, 
    \[
        \lim_{k \to \infty} a_{n_k} = L
    \]
\end{theorem}

\begin{proof}
    Fix $\ep > 0$, we seek $K > 0$ s.t. $ \left| a_{n_k}  - L \right| < \ep$,
    $\forall k > K$. 
    Since $a_n$ converges, 
    \[
        \lim_{n \to \infty} a_n = L \implies \exists N > 0 \text{ s.t }
        |a_n - L | < \ep, \forall n > N
    \]
    Also, $\forall k > N, n_k \geq k \geq N$, hence $|a_{n_k} - L | < \ep$, 
    hence 
    \[
        \lim_{k \to \infty} a_{n_k = L}
    \]
\end{proof}

\begin{corollary}[Convergent subsequences whose limits disagree imply divergence]
    If $\{ a_n \} $ has two subsequences that converge to different values
    then $\{ a_n \} $ diverges and does not converge to $\pm \infty$.
\end{corollary}

\begin{theorem}[Bdd divergent sequences have at least two convergent subsequences]
    \label{theroem:bdd-dvgt-implies-multiple-cvgt-subseq}
    If $\{ a_n \} $ bounded and divergent, 
    then there exist two subsequences of $\{ a_n \} $ that converge 
    to different values. 
\end{theorem}

\begin{theorem}
    \label{theorem:bdd-implies-cvgt-subseq}
    Every bounded sequence $\{ a_n \} $ has a convergent subsequence.
\end{theorem}













